\documentclass[10pt,a4paper]{amsart}
\usepackage[margin=19mm]{geometry}
\usepackage{amsmath,amssymb,mathtools}
\usepackage[scr=rsfs,cal=euler]{mathalfa}
\usepackage{xfrac}
\usepackage[unicode,hidelinks,bookmarks=false]{hyperref}
\newcommand{\ie}{i.e.\ }
\newcommand{\cf}{cf.\ }
\newcommand{\resp}{respectively\ }
\newcommand{\Subsection}{\subsection}
\providecommand{\nobreakdash}{\nobreak\hbox{-}}
\setlength{\emergencystretch}{2em}
\multlinegap0pt
\usepackage{amssymb}
\usepackage{mathtools}
\usepackage{xfrac}
\usepackage{tikz}
\usepackage[shortlabels]{enumitem}
\newtheorem{lemma}{Lemma}
\title{$\mathbb T$-homogeneous locally nilpotent derivations of trinomial algebras}
\author{Timofey Krasikov, Kirill Rassolov}
\date{Source: arXiv:2605.17670v2 (CC BY 4.0)}
\begin{document}
\maketitle

\noindent\emph{Source and licence.} $\mathbb T$-homogeneous locally nilpotent derivations of trinomial algebras. Version arXiv:2605.17670v2 (CC BY 4.0), distributed by arXiv under the Creative Commons Attribution 4.0 International licence, http://creativecommons.org/licenses/by/4.0/, as recorded in the arXiv licence metadata for this exact version and shown on the abstract page https://arxiv.org/abs/2605.17670v2. Full licence notice: \texttt{licenses/arxiv-2605.17670-CC-BY-4.0.txt}.\par\smallskip

\noindent\textit{Editorial scope.} Counted unit: the article's lemma LemmaQuasiCoprime (source0.tex lines 810--812). The article's complete original proof of it is delivered with it (source0.tex lines 813--820, one \texttt{proof} environment of 8 lines). No mathematics is written, repaired or re-derived: the statement and the proof are the article's own lines, in the article's own notation, and no retained line is substituted or re-flowed.  No further source-local context is needed: every cross-reference of every retained line resolves inside the two delivered spans.  Nothing was omitted from any retained span, so the package carries no omission record.  No retained line contains a citation, so the wrapper carries no bibliography and no bibliography entry.  No retained line contains a figure environment, an image-inclusion command or drawing code, so no image is delivered and the package declares no asset dependency.  Editorial formatting only, in addition: an AMS \texttt{amsart} preamble that restates the article's own declarations and macros used by the retained lines, namely the environments \texttt{lemma} on separate counters, together with the standard packages those lines need.  The retained environments print in this document as Lemma 1 in order of appearance, whereas the article prints the same items with the numbering of its own full document.  The complete native source of the article as distributed in the arXiv e-print for this exact version is bundled unchanged as \texttt{sources/200810/source0.tex}.
    \begin{lemma}\label{LemmaQuasiCoprime}
        Let a monomial $P$ and polynomials $Q, F$ be such that $P\mid T_2^{l_2}$, $F$ contains no monomials divisible by~$T_2^{l_2}/P$, and $Q$ does not depend on any of the variables $T_{2j}$.  Suppose that $FP=hQ$ in $R$ for some $h\in R$. Then $P\mid h$ in $R$.
        \end{lemma}

        \begin{proof}
            Choose a polynomial $H$ representing $h\in R$ in such a way that it contains no monomials divisible by $T_2^{l_2}$. Then
            $
           FP-HQ\in I
            $
            and each term in $FP-HQ$ contains no monomials divisible by $T_2^{l_2}$. Consequently,
            $FP=HQ.$ Since the polynomials $P$ and $Q$ are coprime, we obtain $P\mid H$.
        \end{proof}

\end{document}
